\par\noindent\textbf{(1)} 设 u=1/r，中心力 F(r)=−k/r²+2ε/r³。Binet式 u''+u=−mF(1/u)/(ℓ²u²)，故 u''+[1+2mε/ℓ²]u=mk/ℓ²。定义 α²=1+2mε/ℓ²，方程为 u''+α²u=mk/ℓ²。
\par\noindent\textbf{(2)} 解得 u(φ)=mk/(ℓ²α²)[1+e\_* cos(α(φ−φ\_0))]。径向从一次近心点到下一次需 cos 相位增加 2π，因此拱线角 Δφ=2π/α；轨道只在 α 有理相关时闭合，普通Kepler ε=0时 α=1。
\par\noindent\textbf{(3)} 每周相对 2π 的近心进动为 δ=2π(α\textasciicircum{}\{-1\}−1)。若 |2mε/ℓ²|≪1，则 α\textasciicircum{}\{-1\}≈1−mε/ℓ²，故 δ≈−2πmε/ℓ²。正 ε 导致逆向（退行）进动，负 ε 则顺向；此一阶项需在固定角动量下解释。
\par\noindent\textbf{(4)} 有效径向势 U\_eff=ℓ²/(2mr²)+ε/r²−k/r，存在稳定圆轨道仅当 ℓ²/(2m)+ε>0，等价 α²>0。ε→0 恢复封闭椭圆；若 ε≤−ℓ²/(2m)，离心势无法阻止向心坠落，角向振动根不再为实，不能外推微扰进动式。
